\section{十五世纪的区域比较与制度实践}

同一项制度在不同地区从不以同一速度、同一成本或同一社会后果展开。账本中的诏令、军报、赋役记录、工程奏报与使节文书提供跨区域比较的起点；地方志、契约、碑刻、遗址与异域材料则提醒我们，中央分类无法覆盖地方实践。本节以事件链而非单一结论呈现制度、文化、环境与边地关系的差异。

\subsection{漕区、军户与财政实践}

\eventanchor{ML-1450-003}{明·1450（景泰元年）}
\paragraph{1450（景泰元年）：山东发生饥荒} 这一事件使区域差异进入可核查的编年。账本确认的范围是编年候选的年序可由官修与后出材料交叉；具体月日、参与者、战果或数字必须回查所列材料，不能将单一叙事当作定论。，但各材料服务于不同的行政、叙事或保存目的。事件之所以在该时点发生，与；可见的后续变化是赈济、迁徙和损失的实际范围仍须以受灾地区材料分别核验。比较不同地区时，不能把中心政策、地方报数、物质遗存与社会经验混为同一种证据。译写候选以编年材料定位；实录记载反映朝廷选择，崇祯期须另以奏疏、地方及敌对政权材料交叉。

\eventanchor{ML-1450-004}{明·1450（景泰元年）}
\paragraph{1450（景泰元年）：军饷、漕运和户籍赋役的本年记录揭示恢复秩序的行政成本，账面额与实收须分列} 这一事件使区域差异进入可核查的编年。账本确认的范围是来源导向的年度桥接条目：本年材料可定位相关政务线索；具体月日、发文机关、地域和执行结果仍须逐项回查，不能以制度文本或奏报替代实际效果。，但各材料服务于不同的行政、叙事或保存目的。事件之所以在该时点发生，与；可见的后续变化是其法定安排不等于地方执行，后续须以地域材料追踪负担与成效。比较不同地区时，不能把中心政策、地方报数、物质遗存与社会经验混为同一种证据。桥接条目只标示可回查的年度问题与材料链；实录有中央选择性，崇祯期尤其需要奏疏、地方、朝鲜及满文材料交叉。

\eventanchor{ML-1451-001}{明·1451（景泰二年）}
\paragraph{1451（景泰二年）：土木危机后的边镇整顿、京师防务或复辟前后军政调度在本年材料中构成连续问题。（材料核查重点：诏令或奏议的形成与发受链）} 这一事件使区域差异进入可核查的编年。账本确认的范围是来源导向的年度桥接条目：本年材料可定位相关政务线索；具体月日、发文机关、地域和执行结果仍须逐项回查，不能以制度文本或奏报替代实际效果。，但各材料服务于不同的行政、叙事或保存目的。事件之所以在该时点发生，与；可见的后续变化是它改变了后续调兵、补给或地方秩序的条件；战果和伤亡须分开核对。比较不同地区时，不能把中心政策、地方报数、物质遗存与社会经验混为同一种证据。桥接条目只标示可回查的年度问题与材料链；实录有中央选择性，崇祯期尤其需要奏疏、地方、朝鲜及满文材料交叉。；本条特别核对诏令或奏议的形成与发受链。

\eventanchor{ML-1451-002}{明·1451（景泰二年）}
\paragraph{1451（景泰二年）：科举、学校和法律条文在本年制度材料中可追踪，不能据规范文本推定州县实践一致} 这一事件使区域差异进入可核查的编年。账本确认的范围是来源导向的年度桥接条目：本年材料可定位相关政务线索；具体月日、发文机关、地域和执行结果仍须逐项回查，不能以制度文本或奏报替代实际效果。，但各材料服务于不同的行政、叙事或保存目的。事件之所以在该时点发生，与；可见的后续变化是它提供制度和传播史的锚点，不能据此推断社会整体接受。比较不同地区时，不能把中心政策、地方报数、物质遗存与社会经验混为同一种证据。桥接条目只标示可回查的年度问题与材料链；实录有中央选择性，崇祯期尤其需要奏疏、地方、朝鲜及满文材料交叉。

\eventanchor{ML-1451-003}{明·1451（景泰二年）}
\paragraph{1451（景泰二年）：地方灾异、赈济与水利条目为本年社会史提供入口，地方志的修纂层次需要说明} 这一事件使区域差异进入可核查的编年。账本确认的范围是来源导向的年度桥接条目：本年材料可定位相关政务线索；具体月日、发文机关、地域和执行结果仍须逐项回查，不能以制度文本或奏报替代实际效果。，但各材料服务于不同的行政、叙事或保存目的。事件之所以在该时点发生，与；可见的后续变化是赈济、迁徙和损失的实际范围仍须以受灾地区材料分别核验。比较不同地区时，不能把中心政策、地方报数、物质遗存与社会经验混为同一种证据。桥接条目只标示可回查的年度问题与材料链；实录有中央选择性，崇祯期尤其需要奏疏、地方、朝鲜及满文材料交叉。

\eventanchor{ML-1451-004}{明·1451（景泰二年）}
\paragraph{1451（景泰二年）：土木危机后的边镇整顿、京师防务或复辟前后军政调度在本年材料中构成连续问题。（材料核查重点：报称信息的来源、抄传与行政分类）} 这一事件使区域差异进入可核查的编年。账本确认的范围是来源导向的年度桥接条目：本年材料可定位相关政务线索；具体月日、发文机关、地域和执行结果仍须逐项回查，不能以制度文本或奏报替代实际效果。，但各材料服务于不同的行政、叙事或保存目的。事件之所以在该时点发生，与；可见的后续变化是它改变了后续调兵、补给或地方秩序的条件；战果和伤亡须分开核对。比较不同地区时，不能把中心政策、地方报数、物质遗存与社会经验混为同一种证据。桥接条目只标示可回查的年度问题与材料链；实录有中央选择性，崇祯期尤其需要奏疏、地方、朝鲜及满文材料交叉。；本条特别核对报称信息的来源、抄传与行政分类。

\subsection{北边、西南与边地中介}

\eventanchor{ML-1451-005}{明·1451（景泰二年）}
\paragraph{1451（景泰二年）：蒙古诸部、朝鲜及西北绿洲的交涉在本年留下多方材料，应以具体使团和地名校正概称} 这一事件使区域差异进入可核查的编年。账本确认的范围是来源导向的年度桥接条目：本年材料可定位相关政务线索；具体月日、发文机关、地域和执行结果仍须逐项回查，不能以制度文本或奏报替代实际效果。，但各材料服务于不同的行政、叙事或保存目的。事件之所以在该时点发生，与；可见的后续变化是它为后续册封、互市或海防安排留下文书与跨政权比较线索。比较不同地区时，不能把中心政策、地方报数、物质遗存与社会经验混为同一种证据。桥接条目只标示可回查的年度问题与材料链；实录有中央选择性，崇祯期尤其需要奏疏、地方、朝鲜及满文材料交叉。

\eventanchor{ML-1452-001}{明·1452（景泰三年）}
\paragraph{1452（景泰三年）：1452年黄河洪水：黄河堤决口} 这一事件使区域差异进入可核查的编年。账本确认的范围是编年候选的年序可由官修与后出材料交叉；具体月日、参与者、战果或数字必须回查所列材料，不能将单一叙事当作定论。，但各材料服务于不同的行政、叙事或保存目的。事件之所以在该时点发生，与；可见的后续变化是赈济、迁徙和损失的实际范围仍须以受灾地区材料分别核验。比较不同地区时，不能把中心政策、地方报数、物质遗存与社会经验混为同一种证据。译写候选以编年材料定位；实录记载反映朝廷选择，崇祯期须另以奏疏、地方及敌对政权材料交叉。

\eventanchor{ML-1452-002}{明·1452（景泰三年）}
\paragraph{1452（景泰三年）：瑶族、苗族叛乱被镇压} 这一事件使区域差异进入可核查的编年。账本确认的范围是编年候选的年序可由官修与后出材料交叉；具体月日、参与者、战果或数字必须回查所列材料，不能将单一叙事当作定论。，但各材料服务于不同的行政、叙事或保存目的。事件之所以在该时点发生，与；可见的后续变化是它影响后续人事与政策议程；动机和派系标签不能由单一叙事裁定。比较不同地区时，不能把中心政策、地方报数、物质遗存与社会经验混为同一种证据。译写候选以编年材料定位；实录记载反映朝廷选择，崇祯期须另以奏疏、地方及敌对政权材料交叉。

\eventanchor{ML-1452-003}{明·1452（景泰三年）}
\paragraph{1452（景泰三年）：中国北方遭遇洪涝灾害} 这一事件使区域差异进入可核查的编年。账本确认的范围是编年候选的年序可由官修与后出材料交叉；具体月日、参与者、战果或数字必须回查所列材料，不能将单一叙事当作定论。，但各材料服务于不同的行政、叙事或保存目的。事件之所以在该时点发生，与；可见的后续变化是赈济、迁徙和损失的实际范围仍须以受灾地区材料分别核验。比较不同地区时，不能把中心政策、地方报数、物质遗存与社会经验混为同一种证据。译写候选以编年材料定位；实录记载反映朝廷选择，崇祯期须另以奏疏、地方及敌对政权材料交叉。

\eventanchor{ML-1452-004}{明·1452（景泰三年）}
\paragraph{1452（景泰三年）：科举、学校和法律条文在本年制度材料中可追踪，不能据规范文本推定州县实践一致} 这一事件使区域差异进入可核查的编年。账本确认的范围是来源导向的年度桥接条目：本年材料可定位相关政务线索；具体月日、发文机关、地域和执行结果仍须逐项回查，不能以制度文本或奏报替代实际效果。，但各材料服务于不同的行政、叙事或保存目的。事件之所以在该时点发生，与；可见的后续变化是它提供制度和传播史的锚点，不能据此推断社会整体接受。比较不同地区时，不能把中心政策、地方报数、物质遗存与社会经验混为同一种证据。桥接条目只标示可回查的年度问题与材料链；实录有中央选择性，崇祯期尤其需要奏疏、地方、朝鲜及满文材料交叉。

\eventanchor{MM-1452-HEIR-CHANGE}{明·1452（景泰三年）五月}
\paragraph{1452（景泰三年）五月：景泰帝废朱见深太子位，改立其子朱见济为皇太子} 这一事件使区域差异进入可核查的编年。账本确认的范围是废立诏令与日期可靠；促成废立的宫廷协商及各派态度不能仅凭复辟后叙事判定。，但各材料服务于不同的行政、叙事或保存目的。事件之所以在该时点发生，与；可见的后续变化是废立加深了景泰朝与英宗一系的继承对立；朱见济早夭后，储位问题重新悬置。比较不同地区时，不能把中心政策、地方报数、物质遗存与社会经验混为同一种证据。同一名称下的区域执行、数字口径与社会后果仍须分别查证。

\subsection{城市、道路与文化资源}

\eventanchor{ML-1453-001}{明·1453（景泰四年）}
\paragraph{1453（景泰四年）：地方灾异、赈济与水利条目为本年社会史提供入口，地方志的修纂层次需要说明。（材料核查重点：诏令或奏议的形成与发受链）} 这一事件使区域差异进入可核查的编年。账本确认的范围是来源导向的年度桥接条目：本年材料可定位相关政务线索；具体月日、发文机关、地域和执行结果仍须逐项回查，不能以制度文本或奏报替代实际效果。，但各材料服务于不同的行政、叙事或保存目的。事件之所以在该时点发生，与；可见的后续变化是赈济、迁徙和损失的实际范围仍须以受灾地区材料分别核验。比较不同地区时，不能把中心政策、地方报数、物质遗存与社会经验混为同一种证据。桥接条目只标示可回查的年度问题与材料链；实录有中央选择性，崇祯期尤其需要奏疏、地方、朝鲜及满文材料交叉。；本条特别核对诏令或奏议的形成与发受链。

\eventanchor{ML-1453-002}{明·1453（景泰四年）}
\paragraph{1453（景泰四年）：土木危机后的边镇整顿、京师防务或复辟前后军政调度在本年材料中构成连续问题} 这一事件使区域差异进入可核查的编年。账本确认的范围是来源导向的年度桥接条目：本年材料可定位相关政务线索；具体月日、发文机关、地域和执行结果仍须逐项回查，不能以制度文本或奏报替代实际效果。，但各材料服务于不同的行政、叙事或保存目的。事件之所以在该时点发生，与；可见的后续变化是它改变了后续调兵、补给或地方秩序的条件；战果和伤亡须分开核对。比较不同地区时，不能把中心政策、地方报数、物质遗存与社会经验混为同一种证据。桥接条目只标示可回查的年度问题与材料链；实录有中央选择性，崇祯期尤其需要奏疏、地方、朝鲜及满文材料交叉。

\eventanchor{ML-1453-003}{明·1453（景泰四年）}
\paragraph{1453（景泰四年）：景泰、天顺时期的继承、人事与诏令链条可在本年材料中定位，政变责任不可由单一史书裁断} 这一事件使区域差异进入可核查的编年。账本确认的范围是来源导向的年度桥接条目：本年材料可定位相关政务线索；具体月日、发文机关、地域和执行结果仍须逐项回查，不能以制度文本或奏报替代实际效果。，但各材料服务于不同的行政、叙事或保存目的。事件之所以在该时点发生，与；可见的后续变化是它影响后续人事与政策议程；动机和派系标签不能由单一叙事裁定。比较不同地区时，不能把中心政策、地方报数、物质遗存与社会经验混为同一种证据。桥接条目只标示可回查的年度问题与材料链；实录有中央选择性，崇祯期尤其需要奏疏、地方、朝鲜及满文材料交叉。

\eventanchor{ML-1453-004}{明·1453（景泰四年）}
\paragraph{1453（景泰四年）：地方灾异、赈济与水利条目为本年社会史提供入口，地方志的修纂层次需要说明。（材料核查重点：报称信息的来源、抄传与行政分类）} 这一事件使区域差异进入可核查的编年。账本确认的范围是来源导向的年度桥接条目：本年材料可定位相关政务线索；具体月日、发文机关、地域和执行结果仍须逐项回查，不能以制度文本或奏报替代实际效果。，但各材料服务于不同的行政、叙事或保存目的。事件之所以在该时点发生，与；可见的后续变化是赈济、迁徙和损失的实际范围仍须以受灾地区材料分别核验。比较不同地区时，不能把中心政策、地方报数、物质遗存与社会经验混为同一种证据。桥接条目只标示可回查的年度问题与材料链；实录有中央选择性，崇祯期尤其需要奏疏、地方、朝鲜及满文材料交叉。；本条特别核对报称信息的来源、抄传与行政分类。

\eventanchor{MM-1453-ESEN-KHAN}{明·1453（景泰四年）}
\paragraph{1453（景泰四年）：也先杀脱脱不花后自称大汗，瓦剌内部权力关系重组} 这一事件使区域差异进入可核查的编年。账本确认的范围是脱脱不花遇害及也先称汗见于明、蒙古叙事；具体年月和草原诸部的服从范围存在异文。，但各材料服务于不同的行政、叙事或保存目的。事件之所以在该时点发生，与；可见的后续变化是草原权力斗争削弱了一个统一的瓦剌指挥中心，也使明廷面对的北方交涉对象更不稳定。比较不同地区时，不能把中心政策、地方报数、物质遗存与社会经验混为同一种证据。同一名称下的区域执行、数字口径与社会后果仍须分别查证。

\eventanchor{ML-1454-001}{明·1454（景泰五年）}
\paragraph{1454（景泰五年）：异常大雪导致苏州、杭州饥荒} 这一事件使区域差异进入可核查的编年。账本确认的范围是编年候选的年序可由官修与后出材料交叉；具体月日、参与者、战果或数字必须回查所列材料，不能将单一叙事当作定论。，但各材料服务于不同的行政、叙事或保存目的。事件之所以在该时点发生，与；可见的后续变化是它影响后续人事与政策议程；动机和派系标签不能由单一叙事裁定。比较不同地区时，不能把中心政策、地方报数、物质遗存与社会经验混为同一种证据。译写候选以编年材料定位；实录记载反映朝廷选择，崇祯期须另以奏疏、地方及敌对政权材料交叉。

\subsection{环境、工程与地方应对}

\eventanchor{ML-1454-002}{明·1454（景泰五年）}
\paragraph{1454（景泰五年）：土木危机后的边镇整顿、京师防务或复辟前后军政调度在本年材料中构成连续问题} 这一事件使区域差异进入可核查的编年。账本确认的范围是来源导向的年度桥接条目：本年材料可定位相关政务线索；具体月日、发文机关、地域和执行结果仍须逐项回查，不能以制度文本或奏报替代实际效果。，但各材料服务于不同的行政、叙事或保存目的。事件之所以在该时点发生，与；可见的后续变化是它改变了后续调兵、补给或地方秩序的条件；战果和伤亡须分开核对。比较不同地区时，不能把中心政策、地方报数、物质遗存与社会经验混为同一种证据。桥接条目只标示可回查的年度问题与材料链；实录有中央选择性，崇祯期尤其需要奏疏、地方、朝鲜及满文材料交叉。

\eventanchor{ML-1454-003}{明·1454（景泰五年）}
\paragraph{1454（景泰五年）：景泰、天顺时期的继承、人事与诏令链条可在本年材料中定位，政变责任不可由单一史书裁断} 这一事件使区域差异进入可核查的编年。账本确认的范围是来源导向的年度桥接条目：本年材料可定位相关政务线索；具体月日、发文机关、地域和执行结果仍须逐项回查，不能以制度文本或奏报替代实际效果。，但各材料服务于不同的行政、叙事或保存目的。事件之所以在该时点发生，与；可见的后续变化是它影响后续人事与政策议程；动机和派系标签不能由单一叙事裁定。比较不同地区时，不能把中心政策、地方报数、物质遗存与社会经验混为同一种证据。桥接条目只标示可回查的年度问题与材料链；实录有中央选择性，崇祯期尤其需要奏疏、地方、朝鲜及满文材料交叉。

\eventanchor{ML-1454-004}{明·1454（景泰五年）}
\paragraph{1454（景泰五年）：军饷、漕运和户籍赋役的本年记录揭示恢复秩序的行政成本，账面额与实收须分列} 这一事件使区域差异进入可核查的编年。账本确认的范围是来源导向的年度桥接条目：本年材料可定位相关政务线索；具体月日、发文机关、地域和执行结果仍须逐项回查，不能以制度文本或奏报替代实际效果。，但各材料服务于不同的行政、叙事或保存目的。事件之所以在该时点发生，与；可见的后续变化是其法定安排不等于地方执行，后续须以地域材料追踪负担与成效。比较不同地区时，不能把中心政策、地方报数、物质遗存与社会经验混为同一种证据。桥接条目只标示可回查的年度问题与材料链；实录有中央选择性，崇祯期尤其需要奏疏、地方、朝鲜及满文材料交叉。

\eventanchor{ML-1455-001}{明·1455（景泰六年）}
\paragraph{1455（景泰六年）：徐有贞完成黄河堤防修缮工作} 这一事件使区域差异进入可核查的编年。账本确认的范围是编年候选的年序可由官修与后出材料交叉；具体月日、参与者、战果或数字必须回查所列材料，不能将单一叙事当作定论。，但各材料服务于不同的行政、叙事或保存目的。事件之所以在该时点发生，与；可见的后续变化是其法定安排不等于地方执行，后续须以地域材料追踪负担与成效。比较不同地区时，不能把中心政策、地方报数、物质遗存与社会经验混为同一种证据。译写候选以编年材料定位；实录记载反映朝廷选择，崇祯期须另以奏疏、地方及敌对政权材料交叉。

\eventanchor{ML-1455-002}{明·1455（景泰六年）}
\paragraph{1455（景泰六年）：华中地区大面积干旱} 这一事件使区域差异进入可核查的编年。账本确认的范围是编年候选的年序可由官修与后出材料交叉；具体月日、参与者、战果或数字必须回查所列材料，不能将单一叙事当作定论。，但各材料服务于不同的行政、叙事或保存目的。事件之所以在该时点发生，与；可见的后续变化是赈济、迁徙和损失的实际范围仍须以受灾地区材料分别核验。比较不同地区时，不能把中心政策、地方报数、物质遗存与社会经验混为同一种证据。译写候选以编年材料定位；实录记载反映朝廷选择，崇祯期须另以奏疏、地方及敌对政权材料交叉。

\eventanchor{ML-1455-003}{明·1455（景泰六年）}
\paragraph{1455（景泰六年）：蒙古诸部、朝鲜及西北绿洲的交涉在本年留下多方材料，应以具体使团和地名校正概称} 这一事件使区域差异进入可核查的编年。账本确认的范围是来源导向的年度桥接条目：本年材料可定位相关政务线索；具体月日、发文机关、地域和执行结果仍须逐项回查，不能以制度文本或奏报替代实际效果。，但各材料服务于不同的行政、叙事或保存目的。事件之所以在该时点发生，与；可见的后续变化是它为后续册封、互市或海防安排留下文书与跨政权比较线索。比较不同地区时，不能把中心政策、地方报数、物质遗存与社会经验混为同一种证据。桥接条目只标示可回查的年度问题与材料链；实录有中央选择性，崇祯期尤其需要奏疏、地方、朝鲜及满文材料交叉。

\subsection{环境、工程与地方应对}

\paragraph{1455（景泰六年）：也先在草原内部冲突中被杀，瓦剌汗权继承发生断裂} 这一事件使区域差异进入可核查的编年。账本确认的范围是死亡及大致年份可由明方边报和蒙古汗统叙事互校；杀害者、月日和各部归属在材料中不一。，但各材料服务于不同的行政、叙事或保存目的。事件之所以在该时点发生，与；可见的后续变化是瓦剌的集中领导削弱，明廷北边交涉对象趋于分散；这不等于草原威胁立即消失。比较不同地区时，不能把中心政策、地方报数、物质遗存与社会经验混为同一种证据。同一名称下的区域执行、数字口径与社会后果仍须分别查证。

\paragraph{1456（景泰七年）：湖广苗族叛乱遭镇压} 这一事件使区域差异进入可核查的编年。账本确认的范围是编年候选的年序可由官修与后出材料交叉；具体月日、参与者、战果或数字必须回查所列材料，不能将单一叙事当作定论。，但各材料服务于不同的行政、叙事或保存目的。事件之所以在该时点发生，与；可见的后续变化是它影响后续人事与政策议程；动机和派系标签不能由单一叙事裁定。比较不同地区时，不能把中心政策、地方报数、物质遗存与社会经验混为同一种证据。译写候选以编年材料定位；实录记载反映朝廷选择，崇祯期须另以奏疏、地方及敌对政权材料交叉。

\paragraph{1456（景泰七年）：景泰、天顺时期的继承、人事与诏令链条可在本年材料中定位，政变责任不可由单一史书裁断} 这一事件使区域差异进入可核查的编年。账本确认的范围是来源导向的年度桥接条目：本年材料可定位相关政务线索；具体月日、发文机关、地域和执行结果仍须逐项回查，不能以制度文本或奏报替代实际效果。，但各材料服务于不同的行政、叙事或保存目的。事件之所以在该时点发生，与；可见的后续变化是它影响后续人事与政策议程；动机和派系标签不能由单一叙事裁定。比较不同地区时，不能把中心政策、地方报数、物质遗存与社会经验混为同一种证据。桥接条目只标示可回查的年度问题与材料链；实录有中央选择性，崇祯期尤其需要奏疏、地方、朝鲜及满文材料交叉。

\paragraph{1456（景泰七年）：军饷、漕运和户籍赋役的本年记录揭示恢复秩序的行政成本，账面额与实收须分列} 这一事件使区域差异进入可核查的编年。账本确认的范围是来源导向的年度桥接条目：本年材料可定位相关政务线索；具体月日、发文机关、地域和执行结果仍须逐项回查，不能以制度文本或奏报替代实际效果。，但各材料服务于不同的行政、叙事或保存目的。事件之所以在该时点发生，与；可见的后续变化是其法定安排不等于地方执行，后续须以地域材料追踪负担与成效。比较不同地区时，不能把中心政策、地方报数、物质遗存与社会经验混为同一种证据。桥接条目只标示可回查的年度问题与材料链；实录有中央选择性，崇祯期尤其需要奏疏、地方、朝鲜及满文材料交叉。

\paragraph{1456（景泰七年）：科举、学校和法律条文在本年制度材料中可追踪，不能据规范文本推定州县实践一致} 这一事件使区域差异进入可核查的编年。账本确认的范围是来源导向的年度桥接条目：本年材料可定位相关政务线索；具体月日、发文机关、地域和执行结果仍须逐项回查，不能以制度文本或奏报替代实际效果。，但各材料服务于不同的行政、叙事或保存目的。事件之所以在该时点发生，与；可见的后续变化是它提供制度和传播史的锚点，不能据此推断社会整体接受。比较不同地区时，不能把中心政策、地方报数、物质遗存与社会经验混为同一种证据。桥接条目只标示可回查的年度问题与材料链；实录有中央选择性，崇祯期尤其需要奏疏、地方、朝鲜及满文材料交叉。

\paragraph{1457（天顺元年）二月11日：前皇帝被军方复职，成为天顺皇帝} 这一事件使区域差异进入可核查的编年。账本确认的范围是编年候选的年序可由官修与后出材料交叉；具体月日、参与者、战果或数字必须回查所列材料，不能将单一叙事当作定论。，但各材料服务于不同的行政、叙事或保存目的。事件之所以在该时点发生，与；可见的后续变化是它改变了后续调兵、补给或地方秩序的条件；战果和伤亡须分开核对。比较不同地区时，不能把中心政策、地方报数、物质遗存与社会经验混为同一种证据。译写候选以编年材料定位；实录记载反映朝廷选择，崇祯期须另以奏疏、地方及敌对政权材料交叉。
